\documentclass[
  man,
  floatsintext,
  longtable,
  a4paper,
  nolmodern,
  notxfonts,
  notimes,
  12pt,
  colorlinks=true,linkcolor=blue,citecolor=blue,urlcolor=blue]{apa7}

\usepackage{amsmath}
\usepackage{amssymb}



\usepackage[bidi=default]{babel}
\babelprovide[main,import]{spanish}


% get rid of language-specific shorthands (see #6817):
\let\LanguageShortHands\languageshorthands
\def\languageshorthands#1{}

\RequirePackage{longtable}
\RequirePackage{threeparttablex}

\makeatletter
\renewcommand{\paragraph}{\@startsection{paragraph}{4}{\parindent}%
	{0\baselineskip \@plus 0.2ex \@minus 0.2ex}%
	{-.5em}%
	{\normalfont\normalsize\bfseries\typesectitle}}

\renewcommand{\subparagraph}[1]{\@startsection{subparagraph}{5}{0.5em}%
	{0\baselineskip \@plus 0.2ex \@minus 0.2ex}%
	{-\z@\relax}%
	{\normalfont\normalsize\bfseries\itshape\hspace{\parindent}{#1}\textit{\addperi}}{\relax}}
\makeatother




\usepackage{longtable, booktabs, multirow, multicol, colortbl, hhline, caption, array, float, xpatch}
\usepackage{subcaption}


\renewcommand\thesubfigure{\Alph{subfigure}}
\setcounter{topnumber}{2}
\setcounter{bottomnumber}{2}
\setcounter{totalnumber}{4}
\renewcommand{\topfraction}{0.85}
\renewcommand{\bottomfraction}{0.85}
\renewcommand{\textfraction}{0.15}
\renewcommand{\floatpagefraction}{0.7}

\usepackage{tcolorbox}
\tcbuselibrary{listings,theorems, breakable, skins}
\usepackage{fontawesome5}

\definecolor{quarto-callout-color}{HTML}{909090}
\definecolor{quarto-callout-note-color}{HTML}{0758E5}
\definecolor{quarto-callout-important-color}{HTML}{CC1914}
\definecolor{quarto-callout-warning-color}{HTML}{EB9113}
\definecolor{quarto-callout-tip-color}{HTML}{00A047}
\definecolor{quarto-callout-caution-color}{HTML}{FC5300}
\definecolor{quarto-callout-color-frame}{HTML}{ACACAC}
\definecolor{quarto-callout-note-color-frame}{HTML}{4582EC}
\definecolor{quarto-callout-important-color-frame}{HTML}{D9534F}
\definecolor{quarto-callout-warning-color-frame}{HTML}{F0AD4E}
\definecolor{quarto-callout-tip-color-frame}{HTML}{02B875}
\definecolor{quarto-callout-caution-color-frame}{HTML}{FD7E14}

%\newlength\Oldarrayrulewidth
%\newlength\Oldtabcolsep


\usepackage{hyperref}



\usepackage{color}
\usepackage{fancyvrb}
\newcommand{\VerbBar}{|}
\newcommand{\VERB}{\Verb[commandchars=\\\{\}]}
\DefineVerbatimEnvironment{Highlighting}{Verbatim}{commandchars=\\\{\}}
% Add ',fontsize=\small' for more characters per line
\usepackage{framed}
\definecolor{shadecolor}{RGB}{241,243,245}
\newenvironment{Shaded}{\begin{snugshade}}{\end{snugshade}}
\newcommand{\AlertTok}[1]{\textcolor[rgb]{0.68,0.00,0.00}{#1}}
\newcommand{\AnnotationTok}[1]{\textcolor[rgb]{0.37,0.37,0.37}{#1}}
\newcommand{\AttributeTok}[1]{\textcolor[rgb]{0.40,0.45,0.13}{#1}}
\newcommand{\BaseNTok}[1]{\textcolor[rgb]{0.68,0.00,0.00}{#1}}
\newcommand{\BuiltInTok}[1]{\textcolor[rgb]{0.00,0.23,0.31}{#1}}
\newcommand{\CharTok}[1]{\textcolor[rgb]{0.13,0.47,0.30}{#1}}
\newcommand{\CommentTok}[1]{\textcolor[rgb]{0.37,0.37,0.37}{#1}}
\newcommand{\CommentVarTok}[1]{\textcolor[rgb]{0.37,0.37,0.37}{\textit{#1}}}
\newcommand{\ConstantTok}[1]{\textcolor[rgb]{0.56,0.35,0.01}{#1}}
\newcommand{\ControlFlowTok}[1]{\textcolor[rgb]{0.00,0.23,0.31}{\textbf{#1}}}
\newcommand{\DataTypeTok}[1]{\textcolor[rgb]{0.68,0.00,0.00}{#1}}
\newcommand{\DecValTok}[1]{\textcolor[rgb]{0.68,0.00,0.00}{#1}}
\newcommand{\DocumentationTok}[1]{\textcolor[rgb]{0.37,0.37,0.37}{\textit{#1}}}
\newcommand{\ErrorTok}[1]{\textcolor[rgb]{0.68,0.00,0.00}{#1}}
\newcommand{\ExtensionTok}[1]{\textcolor[rgb]{0.00,0.23,0.31}{#1}}
\newcommand{\FloatTok}[1]{\textcolor[rgb]{0.68,0.00,0.00}{#1}}
\newcommand{\FunctionTok}[1]{\textcolor[rgb]{0.28,0.35,0.67}{#1}}
\newcommand{\ImportTok}[1]{\textcolor[rgb]{0.00,0.46,0.62}{#1}}
\newcommand{\InformationTok}[1]{\textcolor[rgb]{0.37,0.37,0.37}{#1}}
\newcommand{\KeywordTok}[1]{\textcolor[rgb]{0.00,0.23,0.31}{\textbf{#1}}}
\newcommand{\NormalTok}[1]{\textcolor[rgb]{0.00,0.23,0.31}{#1}}
\newcommand{\OperatorTok}[1]{\textcolor[rgb]{0.37,0.37,0.37}{#1}}
\newcommand{\OtherTok}[1]{\textcolor[rgb]{0.00,0.23,0.31}{#1}}
\newcommand{\PreprocessorTok}[1]{\textcolor[rgb]{0.68,0.00,0.00}{#1}}
\newcommand{\RegionMarkerTok}[1]{\textcolor[rgb]{0.00,0.23,0.31}{#1}}
\newcommand{\SpecialCharTok}[1]{\textcolor[rgb]{0.37,0.37,0.37}{#1}}
\newcommand{\SpecialStringTok}[1]{\textcolor[rgb]{0.13,0.47,0.30}{#1}}
\newcommand{\StringTok}[1]{\textcolor[rgb]{0.13,0.47,0.30}{#1}}
\newcommand{\VariableTok}[1]{\textcolor[rgb]{0.07,0.07,0.07}{#1}}
\newcommand{\VerbatimStringTok}[1]{\textcolor[rgb]{0.13,0.47,0.30}{#1}}
\newcommand{\WarningTok}[1]{\textcolor[rgb]{0.37,0.37,0.37}{\textit{#1}}}

\providecommand{\tightlist}{%
  \setlength{\itemsep}{0pt}\setlength{\parskip}{0pt}}
\usepackage{longtable,booktabs,array}
\newcounter{none} % for unnumbered tables
\usepackage{calc} % for calculating minipage widths
% Correct order of tables after \paragraph or \subparagraph
\usepackage{etoolbox}
\makeatletter
\patchcmd\longtable{\par}{\if@noskipsec\mbox{}\fi\par}{}{}
\makeatother
% Allow footnotes in longtable head/foot
\IfFileExists{footnotehyper.sty}{\usepackage{footnotehyper}}{\usepackage{footnote}}
\makesavenoteenv{longtable}

\usepackage{graphicx}
\makeatletter
\newsavebox\pandoc@box
\newcommand*\pandocbounded[1]{% scales image to fit in text height/width
  \sbox\pandoc@box{#1}%
  \Gscale@div\@tempa{\textheight}{\dimexpr\ht\pandoc@box+\dp\pandoc@box\relax}%
  \Gscale@div\@tempb{\linewidth}{\wd\pandoc@box}%
  \ifdim\@tempb\p@<\@tempa\p@\let\@tempa\@tempb\fi% select the smaller of both
  \ifdim\@tempa\p@<\p@\scalebox{\@tempa}{\usebox\pandoc@box}%
  \else\usebox{\pandoc@box}%
  \fi%
}
% Set default figure placement to htbp
\def\fps@figure{htbp}
\makeatother







\usepackage{newtx}

\defaultfontfeatures{Scale=MatchLowercase}
\defaultfontfeatures[\rmfamily]{Ligatures=TeX,Scale=1}





\title{Guía completa de entornos virtuales en R con RStudio (renv)}


\shorttitle{Guía completa de entornos virtuales en R con RStudio (renv)}


\usepackage{etoolbox}



\ccoppy{\textcopyright~2020}



\author{Edison Achalma}



\affiliation{
{Escuela Profesional de Economía, Universidad Nacional de San Cristóbal
de Huamanga}}




\leftheader{Achalma}

\date{2020-06-10}


\abstract{Este abstract será actualizado una vez que se complete el
contenido final del artículo. }

\keywords{keyword1, keyword2}

\authornote{\par{\addORCIDlink{Edison Achalma}{0000-0001-6996-3364}} 
\par{ }
\par{   El autor no tiene conflictos de interés que revelar.    Los
roles de autor se clasificaron utilizando la taxonomía de roles de
colaborador (CRediT; https://credit.niso.org/) de la siguiente
manera:  Edison Achalma:   conceptualización, metodología, análisis
formal, investigación, recursos, curación de
datos, redacción, visualización, supervisión, administración del
proyecto}
\par{La correspondencia relativa a este artículo debe dirigirse a Edison
Achalma, Escuela Profesional de Economía, Universidad Nacional de San
Cristóbal de Huamanga, Portal Independencia N°
57, Ayacucho, AYA 5001, Perú, Email: \href{mailto:elmer.achalma.09@unsch.edu.pe}{elmer.achalma.09@unsch.edu.pe}}
}

\makeatletter
\let\endoldlt\endlongtable
\def\endlongtable{
\hline
\endoldlt
}
\makeatother

\urlstyle{same}



\hypersetup{
  pdftitle={Gu\'{\i}a completa de entornos virtuales en R con RStudio (renv)},
  pdfauthor={Autor}
}
\makeatletter
\@ifpackageloaded{caption}{}{\usepackage{caption}}
\AtBeginDocument{%
\ifdefined\contentsname
  \renewcommand*\contentsname{Tabla de contenidos}
\else
  \newcommand\contentsname{Tabla de contenidos}
\fi
\ifdefined\listfigurename
  \renewcommand*\listfigurename{Lista de Figuras}
\else
  \newcommand\listfigurename{Lista de Figuras}
\fi
\ifdefined\listtablename
  \renewcommand*\listtablename{Lista de Tablas}
\else
  \newcommand\listtablename{Lista de Tablas}
\fi
\ifdefined\figurename
  \renewcommand*\figurename{Figura}
\else
  \newcommand\figurename{Figura}
\fi
\ifdefined\tablename
  \renewcommand*\tablename{Tabla}
\else
  \newcommand\tablename{Tabla}
\fi
}
\@ifpackageloaded{float}{}{\usepackage{float}}
\floatstyle{ruled}
\@ifundefined{c@chapter}{\newfloat{codelisting}{h}{lop}}{\newfloat{codelisting}{h}{lop}[chapter]}
\floatname{codelisting}{Listado}
\newcommand*\listoflistings{\listof{codelisting}{Lista de Listados}}
\makeatother
\makeatletter
\makeatother
\makeatletter
\@ifpackageloaded{caption}{}{\usepackage{caption}}
\@ifpackageloaded{subcaption}{}{\usepackage{subcaption}}
\makeatother
\makeatletter
\@ifpackageloaded{fontawesome5}{}{\usepackage{fontawesome5}}
\makeatother

% From https://tex.stackexchange.com/a/645996/211326
%%% apa7 doesn't want to add appendix section titles in the toc
%%% let's make it do it
\makeatletter
\xpatchcmd{\appendix}
  {\par}
  {\addcontentsline{toc}{section}{\@currentlabelname}\par}
  {}{}
\makeatother

%% Disable longtable counter
%% https://tex.stackexchange.com/a/248395/211326

\usepackage{etoolbox}

\makeatletter
\patchcmd{\LT@caption}
  {\bgroup}
  {\bgroup\global\LTpatch@captiontrue}
  {}{}
\patchcmd{\longtable}
  {\par}
  {\par\global\LTpatch@captionfalse}
  {}{}
\apptocmd{\endlongtable}
  {\ifLTpatch@caption\else\addtocounter{table}{-1}\fi}
  {}{}
\newif\ifLTpatch@caption
\makeatother

\begin{document}

\maketitle


\hypertarget{toc}{}
\tableofcontents
\newpage
\section[Introduction]{Guía completa de entornos virtuales en R con
RStudio (renv)}

\setcounter{secnumdepth}{3}

\setlength\LTleft{0pt}




\section{Guía completa de entornos virtuales en R con
RStudio}\label{guuxeda-completa-de-entornos-virtuales-en-r-con-rstudio}

\begin{quote}
Referencia técnica independiente sobre \textbf{entornos virtuales en R},
centrada en \texttt{renv}, el paquete oficial recomendado por
Posit/RStudio para crear bibliotecas de paquetes aisladas por proyecto.
Basada en la documentación oficial: rstudio.github.io/renv, el
repositorio GitHub rstudio/renv, y la documentación de RStudio IDE sobre
Projects y entornos.
\end{quote}

\begin{center}\rule{0.5\linewidth}{0.5pt}\end{center}

\subsection{Tabla de contenidos}\label{tabla-de-contenidos}

\begin{enumerate}
\def\labelenumi{\arabic{enumi}.}
\tightlist
\item
  \hyperref[1-quuxe9-es-un-entorno-virtual-en-r]{¿Qué es un entorno
  virtual en R?}
\item
  \hyperref[2-renv-el-estuxe1ndar-oficial]{renv: el estándar oficial}
\item
  \hyperref[3-conceptos-clave-biblioteca-repositorio-cachuxe9-lockfile]{Conceptos
  clave: biblioteca, repositorio, caché, lockfile}
\item
  \hyperref[4-instalaciuxf3n-de-renv]{Instalación de renv}
\item
  \hyperref[5-flujo-de-trabajo-buxe1sico]{Flujo de trabajo básico}
\item
  \hyperref[6-integraciuxf3n-con-rstudio-projects]{Integración con
  RStudio Projects}
\item
  \hyperref[7-caso-de-uso-1--tesis-o-trabajo-de-investigaciuxf3n-individual]{Caso
  de uso 1 --- Tesis o trabajo de investigación individual}
\item
  \hyperref[8-caso-de-uso-2--curso-o-asignatura-con-muxfaltiples-pruxe1cticas]{Caso
  de uso 2 --- Curso o asignatura con múltiples prácticas}
\item
  \hyperref[9-caso-de-uso-3--colaboraciuxf3n-en-equipo--proyecto-compartido-en-git]{Caso
  de uso 3 --- Colaboración en equipo / proyecto compartido en Git}
\item
  \hyperref[10-caso-de-uso-4--muxfaltiples-proyectos-en-la-misma-muxe1quina-aislamiento]{Caso
  de uso 4 --- Múltiples proyectos en la misma máquina (aislamiento)}
\item
  \hyperref[11-caso-de-uso-5--reportes-reproducibles-con-quarto--r-markdown]{Caso
  de uso 5 --- Reportes reproducibles con Quarto / R Markdown}
\item
  \hyperref[12-caso-de-uso-6--migrar-un-proyecto-antiguo-sin-renv-a-un-entorno-reproducible]{Caso
  de uso 6 --- Migrar un proyecto antiguo (sin renv) a un entorno
  reproducible}
\item
  \hyperref[13-caso-de-uso-7--perfiles-muxfaltiples-dentro-de-un-mismo-proyecto]{Caso
  de uso 7 --- Perfiles múltiples dentro de un mismo proyecto}
\item
  \hyperref[14-caso-de-uso-8--reproducibilidad-a-largo-plazo-archivar-un-proyecto-terminado]{Caso
  de uso 8 --- Reproducibilidad a largo plazo (archivar un proyecto
  terminado)}
\item
  \hyperref[15-caso-de-uso-9--despliegue-en-servidor--docker--posit-connect]{Caso
  de uso 9 --- Despliegue en servidor / Docker / Posit Connect}
\item
  \hyperref[16-caso-de-uso-10--desarrollo-de-un-paquete-de-r-propio]{Caso
  de uso 10 --- Desarrollo de un paquete de R propio}
\item
  \hyperref[17-referencia-ruxe1pida-de-funciones]{Referencia rápida de
  funciones}
\item
  \hyperref[18-luxedmites-de-renv-lo-que-no-resuelve]{Límites de renv:
  lo que NO resuelve}
\item
  \hyperref[19-soluciuxf3n-de-problemas-comunes]{Solución de problemas
  comunes}
\item
  \hyperref[20-referencias-oficiales]{Referencias oficiales}
\end{enumerate}

\begin{center}\rule{0.5\linewidth}{0.5pt}\end{center}

\subsection{1. ¿Qué es un entorno virtual en
R?}\label{quuxe9-es-un-entorno-virtual-en-r}

En R, todos los paquetes se instalan por defecto en una
\textbf{biblioteca de sistema} (system library) compartida por todas las
sesiones y proyectos. Esto significa que si un proyecto necesita
\texttt{dplyr} versión 1.0 y otro necesita la versión 1.1, ambos se
pisan entre sí: actualizar un paquete para un proyecto puede romper
silenciosamente otro proyecto que dependía de una versión anterior.

Un \textbf{entorno virtual} en R resuelve esto dando a cada proyecto su
propia biblioteca de paquetes, completamente independiente de la
biblioteca de sistema y de la de cualquier otro proyecto. La idea es
análoga a \texttt{venv} en Python o a un \texttt{node\_modules} por
proyecto en Node.js, pero adaptada a las particularidades de R.

El paquete oficial de Posit/RStudio para esto es \textbf{\texttt{renv}}
(abreviatura de \emph{reproducible environment}), sucesor de un paquete
anterior llamado \texttt{packrat}. \texttt{renv} es, en la práctica,
\textbf{el estándar actual} para este propósito en el ecosistema R.

\begin{center}\rule{0.5\linewidth}{0.5pt}\end{center}

\subsection{2. renv: el estándar
oficial}\label{renv-el-estuxe1ndar-oficial}

Según la documentación oficial, \texttt{renv} ayuda a que los proyectos
de R sean:

\begin{itemize}
\tightlist
\item
  \textbf{Aislados (isolated)}: instalar o actualizar un paquete en un
  proyecto no afecta a otros proyectos, porque cada uno tiene su propia
  biblioteca privada.
\item
  \textbf{Portables (portable)}: es fácil transportar un proyecto de una
  máquina a otra, incluso entre distintas plataformas (Linux, Windows,
  etc.).
\item
  \textbf{Reproducibles (reproducible)}: \texttt{renv} registra las
  versiones exactas de los paquetes usados y garantiza que esas mismas
  versiones se instalen en cualquier lugar donde se restaure el
  proyecto.
\end{itemize}

Estos tres pilares ---aislamiento, portabilidad, reproducibilidad--- son
el motivo por el cual \texttt{renv} es relevante en prácticamente
cualquier flujo de trabajo serio en R: desde un script personal hasta un
pipeline de producción.

\begin{center}\rule{0.5\linewidth}{0.5pt}\end{center}

\subsection{3. Conceptos clave: biblioteca, repositorio, caché,
lockfile}\label{conceptos-clave-biblioteca-repositorio-cachuxe9-lockfile}

Antes de usar \texttt{renv} conviene tener claros cuatro términos que la
documentación oficial distingue cuidadosamente:

\subsubsection{\texorpdfstring{3.1. Biblioteca
(\emph{library})}{3.1. Biblioteca (library)}}\label{biblioteca-library}

Una \textbf{biblioteca} es un directorio que contiene paquetes
instalados. Es una fuente común de confusión, porque al escribir
\texttt{library(dplyr)} parece que se está ``cargando la biblioteca
dplyr'', cuando en realidad se está cargando el \emph{paquete} dplyr
desde una biblioteca (carpeta). Por defecto, todos los paquetes van a
una única biblioteca de sistema compartida; con \texttt{renv}, cada
proyecto obtiene su propia \textbf{biblioteca de proyecto}
(\texttt{renv/library}).

Puedes ver tus bibliotecas activas con:

\begin{Shaded}
\begin{Highlighting}[]
\FunctionTok{.libPaths}\NormalTok{()}
\end{Highlighting}
\end{Shaded}

\subsubsection{\texorpdfstring{3.2. Repositorio
(\emph{repository})}{3.2. Repositorio (repository)}}\label{repositorio-repository}

Un \textbf{repositorio} es una fuente de paquetes (por ejemplo CRAN,
Bioconductor, GitHub, o el Posit Public Package Manager).
\texttt{install.packages()} descarga un paquete desde un repositorio y
lo coloca en una biblioteca.

\subsubsection{3.3. Caché global de renv}\label{cachuxe9-global-de-renv}

Para no tener que descargar o compilar el mismo paquete una y otra vez
para cada proyecto, \texttt{renv} mantiene una \textbf{caché global}
compartida en el sistema. Cuando instalas un paquete por primera vez, se
descarga y compila una sola vez; para cada proyecto adicional que lo
necesite, \texttt{renv} simplemente crea un enlace desde la biblioteca
del proyecto hacia esa copia en caché, ahorrando tiempo y espacio en
disco.

\subsubsection{\texorpdfstring{3.4. Lockfile
(\texttt{renv.lock})}{3.4. Lockfile (renv.lock)}}\label{lockfile-renv.lock}

El \textbf{lockfile} es un archivo JSON, siempre llamado
\texttt{renv.lock}, que registra toda la información necesaria para
recrear el proyecto en el futuro: la versión de R usada, los
repositorios de origen, y un registro por cada paquete con su versión
exacta, fuente (CRAN, GitHub, etc.) y hash de verificación. Es el
``contrato'' que permite reproducir el entorno en otra máquina.

\begin{center}\rule{0.5\linewidth}{0.5pt}\end{center}

\subsection{4. Instalación de renv}\label{instalaciuxf3n-de-renv}

\texttt{renv} se instala como cualquier paquete de R, desde la consola
(en RStudio o en R puro, sobre cualquier sistema operativo donde ya
tengas R y RStudio instalados):

\begin{Shaded}
\begin{Highlighting}[]
\FunctionTok{install.packages}\NormalTok{(}\StringTok{"renv"}\NormalTok{)}
\end{Highlighting}
\end{Shaded}

Verifica la instalación:

\begin{Shaded}
\begin{Highlighting}[]
\FunctionTok{library}\NormalTok{(renv)}
\FunctionTok{packageVersion}\NormalTok{(}\StringTok{"renv"}\NormalTok{)}
\end{Highlighting}
\end{Shaded}

No requiere ninguna instalación adicional a nivel de sistema operativo:
a diferencia de \texttt{venv} en Python, \texttt{renv} es un paquete de
R puro y funciona igual en Kubuntu, Arch Linux o Windows, siempre que R
y (opcionalmente) RStudio ya estén instalados.

\begin{center}\rule{0.5\linewidth}{0.5pt}\end{center}

\subsection{5. Flujo de trabajo
básico}\label{flujo-de-trabajo-buxe1sico}

Este es el ciclo de trabajo central documentado oficialmente, y se
repite (con variaciones) en todos los casos de uso de las secciones
siguientes.

\subsubsection{Paso 1 --- Inicializar renv en el
proyecto}\label{paso-1-inicializar-renv-en-el-proyecto}

Desde la consola de R, \textbf{con el directorio de trabajo ubicado en
la carpeta del proyecto}:

\begin{Shaded}
\begin{Highlighting}[]
\NormalTok{renv}\SpecialCharTok{::}\FunctionTok{init}\NormalTok{()}
\end{Highlighting}
\end{Shaded}

Esto hace lo siguiente automáticamente:

\begin{enumerate}
\def\labelenumi{\arabic{enumi}.}
\tightlist
\item
  Crea la biblioteca de proyecto en \texttt{renv/library}.
\item
  Detecta los paquetes que ya estás usando en el proyecto (rastreando
  los archivos \texttt{.R}/\texttt{.Rmd}/\texttt{.qmd} en busca de
  llamadas a \texttt{library()}, \texttt{require()}, etc., mediante
  \texttt{renv::dependencies()}).
\item
  Instala esos paquetes detectados en la biblioteca del proyecto
  (mediante \texttt{renv::hydrate()}, que copia desde tu biblioteca de
  usuario existente cuando es posible, en vez de reinstalar todo desde
  cero).
\item
  Crea el lockfile \texttt{renv.lock} con el estado inicial.
\item
  Crea (o modifica) un \texttt{.Rprofile} local al proyecto, que se
  ejecuta automáticamente cada vez que abres R en esa carpeta, activando
  la biblioteca del proyecto sin que tengas que hacer nada manualmente.
\item
  Reinicia la sesión de R (si estás dentro de RStudio).
\end{enumerate}

\subsubsection{Paso 2 --- Trabajar con
normalidad}\label{paso-2-trabajar-con-normalidad}

Sigues usando \texttt{install.packages()}, \texttt{update.packages()},
etc. como siempre. También puedes usar las variantes propias de renv:

\begin{Shaded}
\begin{Highlighting}[]
\NormalTok{renv}\SpecialCharTok{::}\FunctionTok{install}\NormalTok{(}\StringTok{"dplyr"}\NormalTok{)          }\CommentTok{\# instalar un paquete}
\NormalTok{renv}\SpecialCharTok{::}\FunctionTok{install}\NormalTok{(}\StringTok{"usuario/repo"}\NormalTok{)   }\CommentTok{\# instalar desde GitHub}
\end{Highlighting}
\end{Shaded}

\subsubsection{Paso 3 --- Tomar una ``fotografía'' del entorno
(snapshot)}\label{paso-3-tomar-una-fotografuxeda-del-entorno-snapshot}

Cuando confirmes que tu código funciona con el conjunto de paquetes
actual:

\begin{Shaded}
\begin{Highlighting}[]
\NormalTok{renv}\SpecialCharTok{::}\FunctionTok{snapshot}\NormalTok{()}
\end{Highlighting}
\end{Shaded}

Esto actualiza \texttt{renv.lock} con las versiones exactas de los
paquetes actualmente instalados en la biblioteca del proyecto.

\subsubsection{Paso 4 --- Restaurar el entorno en otra máquina (o más
adelante)}\label{paso-4-restaurar-el-entorno-en-otra-muxe1quina-o-muxe1s-adelante}

Si compartes el proyecto, o lo abres en otra computadora, o simplemente
quieres volver al estado exacto registrado en el lockfile:

\begin{Shaded}
\begin{Highlighting}[]
\NormalTok{renv}\SpecialCharTok{::}\FunctionTok{restore}\NormalTok{()}
\end{Highlighting}
\end{Shaded}

Esto reinstala exactamente las versiones de paquetes anotadas en
\texttt{renv.lock}, sin importar qué versiones tengas instaladas en ese
momento en tu biblioteca de usuario.

\subsubsection{Resumen visual del ciclo}\label{resumen-visual-del-ciclo}

\begin{verbatim}
renv::init()        →  crea infraestructura del proyecto
       ↓
trabajas, instalas/actualizas paquetes
       ↓
renv::snapshot()    →  registra el estado actual en renv.lock
       ↓
(compartes el proyecto / cambias de máquina)
       ↓
renv::restore()     →  reinstala exactamente lo que dice renv.lock
\end{verbatim}

\begin{center}\rule{0.5\linewidth}{0.5pt}\end{center}

\subsection{6. Integración con RStudio
Projects}\label{integraciuxf3n-con-rstudio-projects}

RStudio (el IDE) tiene soporte nativo para \texttt{renv} integrado en el
flujo de \textbf{RStudio Projects}.

Al crear un nuevo proyecto en RStudio (\texttt{File} →
\texttt{New\ Project...}), el asistente incluye una casilla
\textbf{``Use renv with this project''}. Si la marcas, RStudio crea
automáticamente toda la infraestructura de \texttt{renv} (equivalente a
haber corrido \texttt{renv::init()} manualmente) en el momento mismo de
crear el proyecto, antes incluso de escribir una sola línea de código.

Esto es la forma más cómoda de empezar un proyecto nuevo con entorno
virtual desde cero: no hace falta recordar ejecutar ningún comando, el
propio flujo de creación de proyecto en RStudio ya lo deja activado.

Para proyectos ya existentes (creados sin esa casilla marcada),
simplemente ejecuta \texttt{renv::init()} manualmente una vez, como se
describió en la sección 5.

\begin{center}\rule{0.5\linewidth}{0.5pt}\end{center}

\subsection{7. Caso de uso 1 --- Tesis o trabajo de investigación
individual}\label{caso-de-uso-1-tesis-o-trabajo-de-investigaciuxf3n-individual}

\textbf{Escenario}: estás escribiendo tu tesis o un artículo y usas
R/RStudio para el análisis estadístico, con un conjunto de paquetes (por
ejemplo \texttt{tidyverse}, \texttt{lme4}, \texttt{apaTables}) que
pueden actualizarse con el tiempo mientras dura el trabajo (meses o
años).

\textbf{Por qué importa}: si una actualización de un paquete cambia el
comportamiento de una función estadística a mitad de tu tesis, tus
resultados podrían dejar de ser reproducibles exactamente, incluso en tu
propia máquina, varios meses después.

\textbf{Flujo recomendado}:

\begin{Shaded}
\begin{Highlighting}[]
\CommentTok{\# Al iniciar el proyecto de tesis}
\NormalTok{renv}\SpecialCharTok{::}\FunctionTok{init}\NormalTok{()}

\CommentTok{\# Trabajas normalmente, instalando lo que necesites}
\FunctionTok{install.packages}\NormalTok{(}\FunctionTok{c}\NormalTok{(}\StringTok{"tidyverse"}\NormalTok{, }\StringTok{"lme4"}\NormalTok{, }\StringTok{"apaTables"}\NormalTok{))}

\CommentTok{\# Cada vez que termines una etapa importante (p. ej. antes de}
\CommentTok{\# enviar un capítulo a tu asesor, o antes de la sustentación)}
\NormalTok{renv}\SpecialCharTok{::}\FunctionTok{snapshot}\NormalTok{()}
\end{Highlighting}
\end{Shaded}

\textbf{Beneficio concreto}: el día de la sustentación o defensa, si un
jurado o evaluador quiere reproducir exactamente tus resultados en otra
computadora, basta con que reciba la carpeta del proyecto (incluyendo
\texttt{renv.lock}) y ejecute \texttt{renv::restore()}. Obtendrá las
mismas versiones de paquete que usaste tú, eliminando una fuente común
de ``en mi máquina sí funciona''.

\begin{center}\rule{0.5\linewidth}{0.5pt}\end{center}

\subsection{8. Caso de uso 2 --- Curso o asignatura con múltiples
prácticas}\label{caso-de-uso-2-curso-o-asignatura-con-muxfaltiples-pruxe1cticas}

\textbf{Escenario}: dictas o llevas un curso de econometría, estadística
o programación en R con varias prácticas o controles a lo largo del
semestre, cada una con requisitos de paquetes ligeramente distintos (una
práctica usa \texttt{forecast}, otra usa \texttt{plm}, otra usa
\texttt{caret}).

\textbf{Problema sin entornos virtuales}: instalar \texttt{caret} para
la práctica 5 puede arrastrar actualizaciones de dependencias
compartidas que rompan silenciosamente el código de la práctica 2, que
ya habías entregado y calificado.

\textbf{Flujo recomendado}: crear \textbf{un proyecto de RStudio por
práctica o por unidad}, cada uno con su propio \texttt{renv}:

\begin{verbatim}
curso-econometria/
├── practica-01-regresion-lineal/   (renv propio)
├── practica-02-series-tiempo/      (renv propio)
├── practica-03-panel-data/         (renv propio)
└── practica-04-machine-learning/   (renv propio)
\end{verbatim}

En cada subcarpeta, al crearla como proyecto de RStudio, marcas ``Use
renv with this project'', o ejecutas \texttt{renv::init()} manualmente.
Cada práctica queda completamente aislada: actualizar \texttt{caret} en
la práctica 4 no afecta en absoluto a la práctica 1.

\textbf{Beneficio adicional para docentes}: puedes distribuir cada
práctica a tus estudiantes incluyendo el \texttt{renv.lock}
correspondiente, garantizando que todos en el curso usen exactamente las
mismas versiones de paquetes que tú usaste para diseñar y probar los
ejercicios, evitando errores de ``a mí me sale un resultado distinto''
por diferencias de versión.

\begin{center}\rule{0.5\linewidth}{0.5pt}\end{center}

\subsection{9. Caso de uso 3 --- Colaboración en equipo / proyecto
compartido en
Git}\label{caso-de-uso-3-colaboraciuxf3n-en-equipo-proyecto-compartido-en-git}

\textbf{Escenario}: trabajas con coautores o colegas en un repositorio
Git compartido (por ejemplo en GitHub, bajo tu cuenta
\texttt{@achalmed}), con varios contribuyentes ejecutando el mismo
código de análisis en distintas máquinas.

\textbf{Flujo recomendado}:

\begin{enumerate}
\def\labelenumi{\arabic{enumi}.}
\tightlist
\item
  Quien crea el proyecto ejecuta \texttt{renv::init()} y hace el primer
  \texttt{renv::snapshot()}.
\item
  Se hace commit a Git de los archivos generados por renv:
\end{enumerate}

\begin{Shaded}
\begin{Highlighting}[]
\FunctionTok{git}\NormalTok{ add renv.lock .Rprofile renv/settings.json renv/activate.R}
\FunctionTok{git}\NormalTok{ commit }\AttributeTok{{-}m} \StringTok{"Configurar entorno renv del proyecto"}
\FunctionTok{git}\NormalTok{ push}
\end{Highlighting}
\end{Shaded}

\texttt{renv} genera automáticamente un \texttt{.gitignore} apropiado
dentro de la carpeta \texttt{renv/}, de modo que la biblioteca de
paquetes en sí (\texttt{renv/library}, que puede pesar cientos de MB)
\textbf{no} se sube al repositorio; solo se sube el lockfile y la
infraestructura ligera necesaria para reconstruirla.

\begin{enumerate}
\def\labelenumi{\arabic{enumi}.}
\setcounter{enumi}{2}
\item
  Cuando un colaborador clona el repositorio y abre el proyecto en
  RStudio, \texttt{renv} se autoarranca (\emph{bootstrap}): detecta que
  el proyecto usa renv, descarga e instala automáticamente la versión
  apropiada del propio paquete \texttt{renv}, y le pregunta al
  colaborador si desea instalar todos los paquetes necesarios ejecutando
  \texttt{renv::restore()}.
\item
  Cuando alguien instala o actualiza un paquete nuevo para el proyecto,
  debe correr \texttt{renv::snapshot()} y avisar al resto del equipo
  (vía commit/push) que el lockfile cambió, para que los demás corran
  \texttt{renv::restore()} y se sincronicen.
\end{enumerate}

\textbf{Regla práctica para el equipo}: ``si tocas paquetes, haces
snapshot y avisas; si te avisan que el lockfile cambió, haces restore
antes de seguir trabajando''.

\begin{center}\rule{0.5\linewidth}{0.5pt}\end{center}

\subsection{10. Caso de uso 4 --- Múltiples proyectos en la misma
máquina
(aislamiento)}\label{caso-de-uso-4-muxfaltiples-proyectos-en-la-misma-muxe1quina-aislamiento}

\textbf{Escenario}: en tu propia computadora mantienes en paralelo
varios proyectos activos: tus sitios Quarto académicos, scripts de
tutoría en economía, un proyecto de análisis de datos para una
consultoría, y experimentos personales con paquetes nuevos de ciencia de
datos.

\textbf{Problema sin entornos virtuales}: probar la última versión de un
paquete experimental para un proyecto personal puede romper, sin que te
des cuenta, un script de tutoría que dependía de un comportamiento
anterior de ese mismo paquete.

\textbf{Flujo recomendado}: activar \texttt{renv} en \textbf{cada
proyecto independiente} que mantengas activo simultáneamente. Gracias a
la caché global de renv (sección 3.3), esto no implica descargar ni
compilar cada paquete una vez por proyecto: si \texttt{ggplot2} ya está
en caché por haberlo instalado en un proyecto, los demás proyectos
simplemente enlazan a esa misma copia, sin gasto adicional relevante de
espacio en disco ni tiempo de instalación, mientras cada proyecto
conserva su propio registro independiente de qué versión está usando.

\textbf{Resultado}: puedes actualizar libremente los paquetes de tu
proyecto experimental sin ningún riesgo de afectar tus otros proyectos
en producción o entrega.

\begin{center}\rule{0.5\linewidth}{0.5pt}\end{center}

\subsection{11. Caso de uso 5 --- Reportes reproducibles con Quarto / R
Markdown}\label{caso-de-uso-5-reportes-reproducibles-con-quarto-r-markdown}

\textbf{Escenario}: generas documentos con Quarto o R Markdown
(\texttt{.qmd} / \texttt{.Rmd}) que combinan texto y código R, por
ejemplo reportes económicos, material didáctico o capítulos de tesis con
\texttt{apaquarto}.

\textbf{Por qué aplica renv aquí}: \texttt{renv::dependencies()} rastrea
automáticamente los bloques de código dentro de archivos \texttt{.qmd} y
\texttt{.Rmd}, igual que lo hace con scripts \texttt{.R} normales, por
lo que \texttt{renv::init()} y \texttt{renv::hydrate()} detectan
correctamente los paquetes usados dentro de tus documentos
Quarto/RMarkdown sin configuración adicional.

\textbf{Flujo recomendado}:

\begin{Shaded}
\begin{Highlighting}[]
\NormalTok{renv}\SpecialCharTok{::}\FunctionTok{init}\NormalTok{()      }\CommentTok{\# detecta paquetes usados en tus .qmd / .Rmd}
\CommentTok{\# ... renderizas el documento, revisas que todo funcione ...}
\NormalTok{renv}\SpecialCharTok{::}\FunctionTok{snapshot}\NormalTok{()  }\CommentTok{\# fija las versiones que produjeron ese render}
\end{Highlighting}
\end{Shaded}

\textbf{Advertencia importante (limitación documentada oficialmente)}:
\texttt{renv} solo gestiona \textbf{paquetes de R}. El motor de
renderizado \textbf{Pandoc}, del cual depende fuertemente
\texttt{rmarkdown} para convertir el documento a su formato final (PDF,
HTML, Word), \textbf{no} está empaquetado dentro del paquete
\texttt{rmarkdown} ni es gestionado por \texttt{renv}. Esto significa
que restaurar el lockfile garantiza las mismas versiones de paquetes de
R, pero \textbf{no} garantiza, por sí solo, que obtengas exactamente el
mismo render visual si la versión de Pandoc instalada en la máquina de
destino es distinta. Si esto te genera problemas de reproducibilidad
visual, puedes usar el paquete \texttt{pandoc} (gestor de versiones de
Pandoc) como complemento.

\begin{center}\rule{0.5\linewidth}{0.5pt}\end{center}

\subsection{12. Caso de uso 6 --- Migrar un proyecto antiguo (sin renv)
a un entorno
reproducible}\label{caso-de-uso-6-migrar-un-proyecto-antiguo-sin-renv-a-un-entorno-reproducible}

\textbf{Escenario}: tienes un proyecto de R ya avanzado (por ejemplo uno
de tus repositorios \texttt{website-achalma} con scripts de análisis, o
un script de tutoría antiguo) que nunca usó \texttt{renv}, y quieres
empezar a protegerlo hacia adelante sin reescribir nada.

\textbf{Flujo recomendado}:

\begin{Shaded}
\begin{Highlighting}[]
\CommentTok{\# Ubícate (setwd o abre el proyecto) en la carpeta del proyecto existente}
\NormalTok{renv}\SpecialCharTok{::}\FunctionTok{init}\NormalTok{()}
\end{Highlighting}
\end{Shaded}

Como se describió en la sección 5, \texttt{init()} detecta
automáticamente qué paquetes ya está usando el código existente, los
copia (cuando es posible) desde tu biblioteca de usuario actual hacia la
nueva biblioteca de proyecto (evitando reinstalar todo desde cero), y
genera el lockfile con el estado funcional actual del proyecto.

\textbf{Resultado}: a partir de ese momento, ese proyecto queda
``congelado'' en un estado conocido y reproducible, sin haber tenido que
identificar manualmente la lista de dependencias ni sus versiones.

\textbf{Variante --- inicialización en blanco}: si prefieres no instalar
nada automáticamente y declarar tú mismo los paquetes desde cero, puedes
inicializar con biblioteca vacía:

\begin{Shaded}
\begin{Highlighting}[]
\NormalTok{renv}\SpecialCharTok{::}\FunctionTok{init}\NormalTok{(}\AttributeTok{bare =} \ConstantTok{TRUE}\NormalTok{)}
\end{Highlighting}
\end{Shaded}

\begin{center}\rule{0.5\linewidth}{0.5pt}\end{center}

\subsection{13. Caso de uso 7 --- Perfiles múltiples dentro de un mismo
proyecto}\label{caso-de-uso-7-perfiles-muxfaltiples-dentro-de-un-mismo-proyecto}

\textbf{Escenario}: un mismo proyecto necesita comportarse de forma
distinta según el contexto. Ejemplos documentados oficialmente:

\begin{itemize}
\tightlist
\item
  Un perfil de \textbf{``desarrollo''} (\texttt{dev}) mientras pruebas y
  depuras código.
\item
  Un perfil de \textbf{``producción''} (\texttt{prod}) para despliegues
  finales, con un conjunto más controlado de paquetes.
\item
  Un perfil de \textbf{``shiny''} específico si el proyecto también
  contiene una aplicación Shiny con dependencias propias (por ejemplo
  \texttt{shiny}, \texttt{tidyverse}) que no necesitas cargar en el
  resto del proyecto.
\end{itemize}

\textbf{Por qué importa}: en vez de mantener varias copias del proyecto,
\texttt{renv} permite que un mismo proyecto tenga \textbf{múltiples
bibliotecas y lockfiles alternativos}, uno por perfil, todos viviendo
dentro del mismo árbol de carpetas.

\textbf{Flujo recomendado}:

\begin{Shaded}
\begin{Highlighting}[]
\CommentTok{\# Crear y activar el perfil "dev" como perfil por defecto del proyecto}
\NormalTok{renv}\SpecialCharTok{::}\FunctionTok{activate}\NormalTok{(}\AttributeTok{profile =} \StringTok{"dev"}\NormalTok{)}
\end{Highlighting}
\end{Shaded}

Con esto, las rutas de biblioteca y lockfile del proyecto pasan a
resolverse dentro de \texttt{renv/profiles/dev/}, en vez de las rutas
por defecto.

Para activar un perfil solo durante una sesión puntual, sin cambiar el
valor por defecto del proyecto:

\begin{Shaded}
\begin{Highlighting}[]
\FunctionTok{Sys.setenv}\NormalTok{(}\AttributeTok{RENV\_PROFILE =} \StringTok{"dev"}\NormalTok{)}
\end{Highlighting}
\end{Shaded}

o, desde la línea de comandos (por ejemplo en tu \texttt{zsh}), antes de
lanzar R:

\begin{Shaded}
\begin{Highlighting}[]
\BuiltInTok{export} \VariableTok{RENV\_PROFILE}\OperatorTok{=}\NormalTok{dev}
\end{Highlighting}
\end{Shaded}

Para volver al perfil estándar:

\begin{Shaded}
\begin{Highlighting}[]
\NormalTok{renv}\SpecialCharTok{::}\FunctionTok{activate}\NormalTok{(}\AttributeTok{profile =} \StringTok{"default"}\NormalTok{)}
\end{Highlighting}
\end{Shaded}

\textbf{Declarar dependencias específicas de un perfil}: se hace en el
archivo \texttt{DESCRIPTION} del proyecto, con un campo dedicado por
perfil, por ejemplo:

\begin{verbatim}
Config/renv/profiles/shiny/dependencies: shiny, tidyverse
\end{verbatim}

\begin{center}\rule{0.5\linewidth}{0.5pt}\end{center}

\subsection{14. Caso de uso 8 --- Reproducibilidad a largo plazo
(archivar un proyecto
terminado)}\label{caso-de-uso-8-reproducibilidad-a-largo-plazo-archivar-un-proyecto-terminado}

\textbf{Escenario}: terminaste una investigación, tesis o reporte y
quieres archivarlo de forma que, dentro de dos o tres años, alguien
(incluido tú mismo) pueda volver a ejecutarlo exactamente igual, aunque
CRAN ya haya avanzado varias versiones de los paquetes involucrados.

\textbf{Flujo recomendado al cerrar el proyecto}:

\begin{Shaded}
\begin{Highlighting}[]
\NormalTok{renv}\SpecialCharTok{::}\FunctionTok{snapshot}\NormalTok{()   }\CommentTok{\# última fotografía del entorno funcional final}
\end{Highlighting}
\end{Shaded}

Y conservar junto al código fuente:

\begin{itemize}
\tightlist
\item
  El archivo \texttt{renv.lock}.
\item
  El \texttt{.Rprofile} del proyecto.
\item
  La carpeta \texttt{renv/} (sin la subcarpeta \texttt{library}, que es
  regenerable).
\end{itemize}

\textbf{Al reabrir el proyecto archivado, meses o años después}:

\begin{Shaded}
\begin{Highlighting}[]
\NormalTok{renv}\SpecialCharTok{::}\FunctionTok{restore}\NormalTok{()}
\end{Highlighting}
\end{Shaded}

\texttt{renv} reinstalará las versiones exactas registradas, incluso si
en CRAN ya existen versiones más nuevas de esos mismos paquetes, siempre
que los binarios/fuentes originales sigan disponibles en algún
repositorio accesible.

\textbf{Limitación honesta que reconoce la propia documentación}: si un
paquete fue instalado originalmente desde un binario que posteriormente
deja de estar disponible (por ejemplo, retirado de un repositorio),
\texttt{renv} intentará compilarlo desde el código fuente como
alternativa, pero esto puede fallar si faltan las dependencias de
sistema necesarias para compilar (ver la sección de dependencias de
sistema en la guía de instalación de R/RStudio). Por esta razón,
\textbf{renv no es una garantía absoluta de reproducibilidad eterna},
pero sí resuelve la parte más importante y más frecuente del problema:
las versiones de paquetes de R.

\begin{center}\rule{0.5\linewidth}{0.5pt}\end{center}

\subsection{15. Caso de uso 9 --- Despliegue en servidor / Docker /
Posit
Connect}\label{caso-de-uso-9-despliegue-en-servidor-docker-posit-connect}

\textbf{Escenario}: necesitas que tu análisis o aplicación (por ejemplo
una app Shiny, o un pipeline de reportes) se ejecute de forma idéntica
en un servidor, contenedor Docker, integración continua (CI), o en una
plataforma de publicación como Posit Connect.

\textbf{Por qué aplica renv aquí}: el mismo lockfile que usas en tu
laptop es el artefacto que le indica al entorno remoto exactamente qué
instalar. La documentación oficial cubre tres escenarios de despliegue
específicos:

\begin{itemize}
\tightlist
\item
  \textbf{Integración continua (CI)}: en pipelines de CI (por ejemplo
  GitHub Actions), se puede usar \texttt{renv::restore()} como paso de
  configuración del entorno antes de correr pruebas, garantizando que el
  CI use las mismas versiones que tu máquina de desarrollo.
\item
  \textbf{Docker}: se recomienda copiar el \texttt{renv.lock} (y la
  infraestructura de renv) dentro de la imagen y ejecutar
  \texttt{renv::restore()} como parte del \texttt{Dockerfile}, de modo
  que la imagen final tenga exactamente el entorno de paquetes esperado.
  La caché de renv puede aprovecharse entre builds de imágenes para
  acelerar reconstrucciones.
\item
  \textbf{Posit Connect}: al publicar contenido (apps Shiny, reportes
  Quarto/RMarkdown) hacia Posit Connect, el servidor usa el
  \texttt{renv.lock} del proyecto para reconstruir el entorno de
  paquetes necesario en el servidor remoto.
\end{itemize}

\textbf{Flujo conceptual mínimo para cualquiera de estos casos}:

\begin{Shaded}
\begin{Highlighting}[]
\CommentTok{\# En tu máquina de desarrollo, antes de desplegar}
\NormalTok{renv}\SpecialCharTok{::}\FunctionTok{snapshot}\NormalTok{()}
\end{Highlighting}
\end{Shaded}

\begin{Shaded}
\begin{Highlighting}[]
\CommentTok{\# En el entorno remoto / imagen / pipeline}
\CommentTok{\# (usualmente en un script de arranque o en el Dockerfile)}
\end{Highlighting}
\end{Shaded}

\begin{Shaded}
\begin{Highlighting}[]
\NormalTok{renv}\SpecialCharTok{::}\FunctionTok{restore}\NormalTok{()}
\end{Highlighting}
\end{Shaded}

\begin{quote}
Para los detalles específicos de sintaxis de cada plataforma (ejemplo de
\texttt{Dockerfile}, configuración de un workflow de GitHub Actions,
opciones de Posit Connect), consulta los artículos oficiales dedicados:
\emph{Using renv with continuous integration}, \emph{Using renv with
Docker} y \emph{Using renv with Posit Connect}, enlazados en la sección
de referencias.
\end{quote}

\begin{center}\rule{0.5\linewidth}{0.5pt}\end{center}

\subsection{16. Caso de uso 10 --- Desarrollo de un paquete de R
propio}\label{caso-de-uso-10-desarrollo-de-un-paquete-de-r-propio}

\textbf{Escenario}: estás desarrollando tu propio paquete de R (por
ejemplo, una herramienta interna para tus prompts/scripts educativos, o
un paquete que planeas publicar en CRAN).

\textbf{Particularidad de este caso}: cuando \texttt{renv} se usa dentro
de un proyecto de \textbf{desarrollo de paquete} (es decir, un proyecto
que contiene su propio archivo \texttt{DESCRIPTION} de paquete, no solo
de análisis), el comportamiento es ligeramente distinto al de un
proyecto de análisis normal: en vez de usar siempre la biblioteca de
proyecto, \texttt{renv} puede usar una biblioteca externa al directorio
del proyecto, para evitar mezclar la infraestructura del propio paquete
en desarrollo con su entorno de pruebas.

\textbf{Casos típicos dentro de este escenario}:

\begin{itemize}
\tightlist
\item
  \textbf{Probar contra la última versión disponible en CRAN} de tus
  dependencias (recomendado si planeas publicar el paquete en CRAN, ya
  que CRAN exige compatibilidad con el ecosistema actual):
\end{itemize}

\begin{Shaded}
\begin{Highlighting}[]
\NormalTok{renv}\SpecialCharTok{::}\FunctionTok{restore}\NormalTok{()}
\end{Highlighting}
\end{Shaded}

\begin{itemize}
\item
  \textbf{Probar contra versiones de desarrollo} de otras dependencias
  (por ejemplo si dependes de una rama no publicada aún de otro paquete
  tuyo o de un colaborador), declarando esas dependencias de desarrollo
  en el campo \texttt{Remotes} del \texttt{DESCRIPTION} del paquete.
\item
  \textbf{Integración continua del propio paquete}: es habitual usar
  \texttt{renv} dentro del flujo de CI del paquete (por ejemplo en
  GitHub Actions) para acelerar la instalación de dependencias de prueba
  aprovechando la caché global, en vez de reinstalar todo desde cero en
  cada ejecución del pipeline.
\end{itemize}

\begin{quote}
Importante: si vas a enviar tu paquete a CRAN, \textbf{no debe incluirse
la infraestructura de renv} dentro del tarball fuente que se sube a
CRAN; \texttt{renv} es una herramienta para tu flujo de desarrollo, no
parte del paquete distribuido.
\end{quote}

\begin{center}\rule{0.5\linewidth}{0.5pt}\end{center}

\subsection{17. Referencia rápida de
funciones}\label{referencia-ruxe1pida-de-funciones}

{\def\LTcaptype{none} % do not increment counter
\begin{longtable}[]{@{}
  >{\raggedright\arraybackslash}p{(\linewidth - 2\tabcolsep) * \real{0.5000}}
  >{\raggedright\arraybackslash}p{(\linewidth - 2\tabcolsep) * \real{0.5000}}@{}}
\toprule\noalign{}
\begin{minipage}[b]{\linewidth}\raggedright
Función
\end{minipage} & \begin{minipage}[b]{\linewidth}\raggedright
Propósito
\end{minipage} \\
\midrule\noalign{}
\endhead
\bottomrule\noalign{}
\endlastfoot
\texttt{renv::init()} & Inicializa renv en el proyecto actual (crea
biblioteca, lockfile y \texttt{.Rprofile}) \\
\texttt{renv::init(bare\ =\ TRUE)} & Inicializa con biblioteca vacía,
sin instalar nada automáticamente \\
\texttt{renv::snapshot()} & Registra en \texttt{renv.lock} el estado
actual de la biblioteca del proyecto \\
\texttt{renv::restore()} & Reinstala en la biblioteca del proyecto las
versiones exactas del \texttt{renv.lock} \\
\texttt{renv::install("pkg")} & Instala un paquete (alternativa a
\texttt{install.packages()}, soporta GitHub, Bioconductor, etc.) \\
\texttt{renv::update()} & Actualiza todas las dependencias del proyecto
a sus últimas versiones disponibles \\
\texttt{renv::upgrade()} & Actualiza únicamente el propio paquete
\texttt{renv} \\
\texttt{renv::dependencies()} & Lista los paquetes que el código del
proyecto está usando actualmente \\
\texttt{renv::hydrate()} & Instala en la biblioteca de proyecto los
paquetes detectados por \texttt{dependencies()} \\
\texttt{renv::status()} & Compara el estado actual de la biblioteca
contra lo registrado en el lockfile \\
\texttt{renv::history()} & Muestra el historial de versiones anteriores
del lockfile (vía Git) \\
\texttt{renv::revert()} & Revierte el lockfile a una versión anterior
del historial \\
\texttt{renv::activate()} & Activa renv en la sesión actual (o cambia de
perfil con \texttt{profile=}) \\
\texttt{renv::deactivate()} & Desactiva renv en el proyecto (mantiene
los archivos) \\
\texttt{renv::deactivate(clean\ =\ TRUE)} & Elimina por completo la
infraestructura de renv del proyecto \\
\texttt{renv::paths\$root()} & Devuelve la ruta de la caché/raíz global
de renv en el sistema \\
\end{longtable}
}

\begin{center}\rule{0.5\linewidth}{0.5pt}\end{center}

\subsection{18. Límites de renv: lo que NO
resuelve}\label{luxedmites-de-renv-lo-que-no-resuelve}

La propia documentación oficial es explícita sobre esto, y conviene
tenerlo presente para no asumir una falsa sensación de seguridad total:

\begin{itemize}
\tightlist
\item
  \textbf{Versión de R}: \texttt{renv} registra qué versión de R se usó,
  pero no instala ni cambia versiones de R por ti (no puede, porque
  corre dentro de la propia R). Para gestionar múltiples versiones de R
  en una misma máquina, la documentación oficial sugiere herramientas
  externas como \texttt{rig}.
\item
  \textbf{Pandoc}: como se explicó en el caso de uso 5, Pandoc no viaja
  dentro del lockfile de \texttt{renv}, aunque \texttt{rmarkdown}/Quarto
  dependan fuertemente de él.
\item
  \textbf{Sistema operativo, librerías de sistema, versión del
  compilador}: mantener una imagen de sistema ``estable'' es un problema
  aparte; la solución recomendada oficialmente para esto es
  \textbf{Docker}, no \texttt{renv} por sí solo.
\item
  \textbf{Disponibilidad futura de binarios}: si un paquete fue
  instalado desde un binario que luego deja de publicarse,
  \texttt{restore()} intentará compilar desde fuente, lo cual puede
  fallar si faltan dependencias de sistema.
\end{itemize}

En palabras de la propia documentación: hacer un proyecto reproducible
siempre exige reflexión, no solo el uso mecánico de una herramienta.
\texttt{renv} resuelve muy bien \textbf{una} parte importante del
problema (los paquetes de R), no todo el problema de reproducibilidad de
punta a punta.

\begin{center}\rule{0.5\linewidth}{0.5pt}\end{center}

\subsection{19. Solución de problemas
comunes}\label{soluciuxf3n-de-problemas-comunes}

\subsubsection{\texorpdfstring{19.1. \texttt{renv::restore()} falla al
intentar compilar un paquete desde
fuente}{19.1. renv::restore() falla al intentar compilar un paquete desde fuente}}\label{renvrestore-falla-al-intentar-compilar-un-paquete-desde-fuente}

Suele deberse a dependencias de sistema faltantes (ver la guía de
instalación de R/RStudio, sección de ``Dependencias de sistema para
paquetes de R''). Instala la librería de sistema indicada en el mensaje
de error y reintenta.

\subsubsection{19.2. Un colaborador abre el proyecto y no ve activarse
renv
automáticamente}\label{un-colaborador-abre-el-proyecto-y-no-ve-activarse-renv-automuxe1ticamente}

Verifica que se haya hecho commit de \texttt{.Rprofile},
\texttt{renv.lock}, \texttt{renv/settings.json} y
\texttt{renv/activate.R} al repositorio Git. Si alguno de estos archivos
quedó fuera del control de versiones (por ejemplo por un
\texttt{.gitignore} mal configurado), el autoarranque de renv no
ocurrirá en la máquina del colaborador.

\subsubsection{19.3. Quiero saber si mi biblioteca actual coincide con
el
lockfile}\label{quiero-saber-si-mi-biblioteca-actual-coincide-con-el-lockfile}

\begin{Shaded}
\begin{Highlighting}[]
\NormalTok{renv}\SpecialCharTok{::}\FunctionTok{status}\NormalTok{()}
\end{Highlighting}
\end{Shaded}

Esto compara la biblioteca actual del proyecto contra lo registrado en
\texttt{renv.lock} y reporta diferencias (paquetes instalados que no
están en el lockfile, o viceversa).

\subsubsection{19.4. Necesito volver a una versión anterior del lockfile
(no solo a la
última)}\label{necesito-volver-a-una-versiuxf3n-anterior-del-lockfile-no-solo-a-la-uxfaltima}

\begin{Shaded}
\begin{Highlighting}[]
\NormalTok{renv}\SpecialCharTok{::}\FunctionTok{history}\NormalTok{()   }\CommentTok{\# ver versiones anteriores registradas en Git}
\NormalTok{renv}\SpecialCharTok{::}\FunctionTok{revert}\NormalTok{()    }\CommentTok{\# revertir a una versión específica}
\end{Highlighting}
\end{Shaded}

\subsubsection{19.5. El proyecto ya no necesita
renv}\label{el-proyecto-ya-no-necesita-renv}

\begin{Shaded}
\begin{Highlighting}[]
\NormalTok{renv}\SpecialCharTok{::}\FunctionTok{deactivate}\NormalTok{()              }\CommentTok{\# desactiva pero conserva los archivos}
\NormalTok{renv}\SpecialCharTok{::}\FunctionTok{deactivate}\NormalTok{(}\AttributeTok{clean =} \ConstantTok{TRUE}\NormalTok{)  }\CommentTok{\# elimina toda la infraestructura de renv del proyecto}
\end{Highlighting}
\end{Shaded}

Para eliminar también la caché global de renv del sistema (afecta a
\emph{todos} tus proyectos con renv, úsalo con cuidado):

\begin{Shaded}
\begin{Highlighting}[]
\NormalTok{root }\OtherTok{\textless{}{-}}\NormalTok{ renv}\SpecialCharTok{::}\NormalTok{paths}\SpecialCharTok{$}\FunctionTok{root}\NormalTok{()}
\FunctionTok{unlink}\NormalTok{(root, }\AttributeTok{recursive =} \ConstantTok{TRUE}\NormalTok{)}
\end{Highlighting}
\end{Shaded}

Y finalmente, si ya no quieres usar renv en ningún proyecto:

\begin{Shaded}
\begin{Highlighting}[]
\NormalTok{utils}\SpecialCharTok{::}\FunctionTok{remove.packages}\NormalTok{(}\StringTok{"renv"}\NormalTok{)}
\end{Highlighting}
\end{Shaded}

\begin{center}\rule{0.5\linewidth}{0.5pt}\end{center}

\subsection{20. Referencias oficiales}\label{referencias-oficiales}

\begin{itemize}
\tightlist
\item
  \textbf{renv --- Sitio oficial de documentación}:
  https://rstudio.github.io/renv/
\item
  \textbf{renv --- Introducción (vignette principal)}:
  https://rstudio.github.io/renv/articles/renv.html
\item
  \textbf{renv --- Instalación de paquetes y caché}:
  https://rstudio.github.io/renv/articles/package-install.html
\item
  \textbf{renv --- Fuentes de paquetes (CRAN, GitHub, Bioconductor,
  etc.)}: https://rstudio.github.io/renv/articles/package-sources.html
\item
  \textbf{renv --- Perfiles de proyecto}:
  https://rstudio.github.io/renv/articles/profiles.html
\item
  \textbf{renv --- Desarrollo de paquetes con renv}:
  https://rstudio.github.io/renv/articles/packages.html
\item
  \textbf{renv --- Integración continua (CI)}:
  https://rstudio.github.io/renv/articles/ci.html
\item
  \textbf{renv --- Uso con Docker}:
  https://rstudio.github.io/renv/articles/docker.html
\item
  \textbf{renv --- Uso con Posit Connect}:
  https://rstudio.github.io/renv/articles/rsconnect.html
\item
  \textbf{renv --- Preguntas frecuentes (FAQ)}:
  https://rstudio.github.io/renv/articles/faq.html
\item
  \textbf{renv --- Referencia de funciones}:
  https://rstudio.github.io/renv/reference/index.html
\item
  \textbf{renv --- Repositorio oficial en GitHub}:
  https://github.com/rstudio/renv
\item
  \textbf{RStudio IDE --- Guía de entornos con renv}:
  https://github.com/rstudio/rstudio/blob/main/docs/user/rstudio/ide/guide/environments/r/renv.qmd
\end{itemize}

\begin{center}\rule{0.5\linewidth}{0.5pt}\end{center}

\textbf{Edison Achalma} Economista --- Universidad Nacional de San
Cristóbal de Huamanga (UNSCH) Ayacucho, Perú ORCID:
\href{https://orcid.org/0000-0001-6996-3364}{0000-0001-6996-3364}

\section{Publicaciones Similares}\label{publicaciones-similares}

Si te interesó este artículo, te recomendamos que explores otros blogs y
recursos relacionados que pueden ampliar tus conocimientos. Aquí te dejo
algunas sugerencias:

\begin{enumerate}
\def\labelenumi{\arabic{enumi}.}
\tightlist
\item
  \href{https://numerus-scriptum.netlify.app/r/2020-06-10-011-instalacion-de-r/index.pdf}{\faIcon{file-pdf}}
  \href{https://numerus-scriptum.netlify.app/r/2020-06-10-011-instalacion-de-r}{011
  Instalacion De R}
\item
  \href{https://numerus-scriptum.netlify.app/r/2020-06-10-012-configurar-entorno-virtual-r/index.pdf}{\faIcon{file-pdf}}
  \href{https://numerus-scriptum.netlify.app/r/2020-06-10-012-configurar-entorno-virtual-r}{012
  Configurar Entorno Virtual R}
\item
  \href{https://numerus-scriptum.netlify.app/r/2020-06-10-013-lo-que-debemos-saber-de-r/index.pdf}{\faIcon{file-pdf}}
  \href{https://numerus-scriptum.netlify.app/r/2020-06-10-013-lo-que-debemos-saber-de-r}{013
  Lo Que Debemos Saber De R}
\item
  \href{https://numerus-scriptum.netlify.app/r/2021-04-05-02-manipulacion-de-datos/index.pdf}{\faIcon{file-pdf}}
  \href{https://numerus-scriptum.netlify.app/r/2021-04-05-02-manipulacion-de-datos}{02
  Manipulacion De Datos}
\item
  \href{https://numerus-scriptum.netlify.app/r/2021-04-12-03-visualizacion-de-datos/index.pdf}{\faIcon{file-pdf}}
  \href{https://numerus-scriptum.netlify.app/r/2021-04-12-03-visualizacion-de-datos}{03
  Visualizacion De Datos}
\item
  \href{https://numerus-scriptum.netlify.app/r/2022-11-07-04-modelo-de-machine-learning-i-analisis-exploratorio/index.pdf}{\faIcon{file-pdf}}
  \href{https://numerus-scriptum.netlify.app/r/2022-11-07-04-modelo-de-machine-learning-i-analisis-exploratorio}{04
  Modelo De Machine Learning I Analisis Exploratorio}
\item
  \href{https://numerus-scriptum.netlify.app/r/2022-11-14-05-modelo-de-machine-learning-ii-modelo-de-clasificacion/index.pdf}{\faIcon{file-pdf}}
  \href{https://numerus-scriptum.netlify.app/r/2022-11-14-05-modelo-de-machine-learning-ii-modelo-de-clasificacion}{05
  Modelo De Machine Learning Ii Modelo De Clasificacion}
\item
  \href{https://numerus-scriptum.netlify.app/r/2022-11-21-06-modelo-de-machine-learning-iii-modelo-de-regresion/index.pdf}{\faIcon{file-pdf}}
  \href{https://numerus-scriptum.netlify.app/r/2022-11-21-06-modelo-de-machine-learning-iii-modelo-de-regresion}{06
  Modelo De Machine Learning Iii Modelo De Regresion}
\item
  \href{https://numerus-scriptum.netlify.app/r/2022-11-28-07-modelo-de-machine-learning-iv-tex-mining/index.pdf}{\faIcon{file-pdf}}
  \href{https://numerus-scriptum.netlify.app/r/2022-11-28-07-modelo-de-machine-learning-iv-tex-mining}{07
  Modelo De Machine Learning Iv Tex Mining}
\end{enumerate}

Esperamos que encuentres estas publicaciones igualmente interesantes y
útiles. ¡Disfruta de la lectura!






\end{document}
